\documentclass[main.tex]{subfiles}
\begin{document}

I use the estimated model to show that equilibrium fees and rewards are too high, but that increasing network competition without breaking the competitive bottleneck makes the problem worse.
I simulate five counterfactual policies: capping credit card fees, repealing the Durbin Amendment, merging all three networks, increasing consumer multi-homing through dual-routing mandates, and introducing a central bank digital currency.
All counterfactuals fix consumers' preferences $\beta_i$, baseline income $y$, and merchants' sales benefits to card consumers $\gamma$.
Consumer adoption, merchant acceptance, retail prices, and network prices adjust endogenously.\footnote{
    Although merchants' types $\gamma$ are fixed, their incentives to accept cards can change because the effect of card acceptance on total sales depends on $\gamma$ and the share of consumers using each card.
}
The baseline is the observed 2019 equilibrium, where networks maximize profits subject to the Durbin Amendment's debit fee cap while credit fees remain unregulated.
Welfare comparisons are relative to this status quo.

\subsection{Capping Credit Card Merchant Fees}
\label{subsec:cap-fees-intro}

Capping credit card merchant fees at 120 bp, roughly half their current level, while leaving debit fee caps unchanged corrects the price-coherence adoption externality and raises total welfare by \$27 billion per year.
% HARDCODED[table: welfare_table_baseline.tex, row: "Total", col: 1] current=27
Lower merchant fees reduce credit card rewards, correcting excessive adoption driven by price coherence.
The gains are progressive.
Low-income consumers benefit from lower retail prices, while high-income consumers lose from reduced rewards.

\subsubsection{Effects on~Prices~and~Shares}
\label{subsubsec:cf-price-share-effects}

Capping credit card merchant fees reduces rewards and card use, illustrating the see-saw principle in \textcite{Rochet.Tirole2003}.
% HARDCODED[table: cf_table_baseline.tex, row: "Credit Rewards", col: 1, derived: -110 bps = 1.10 pp] current=-110
Table \ref{tab:cf-effects} shows that networks respond to the credit fee cap by cutting credit card rewards by around 1.1 pp., largely eliminating rewards.
Spending on credit falls by nearly half as consumers substitute toward debit and cash (Table \ref{tab:cf-effects}).
% HARDCODED[table: cf_table_baseline.tex, row: "Credit", col: 1] current=-19
Price and quantity changes together reduce total merchant fees and rewards by roughly equal amounts.
I normalize total model income to the \$10 trillion in 2019 consumer spending \parencite{Nilson2020a}, so each bp corresponds to \$1 billion.
% CONSTANT: normalization from Nilson (2020a), not updated by pipeline
Hence fees and rewards fall by around \$58 billion each.
% HARDCODED[table: cf_table_baseline.tex, row: "Total Fees", col: 1] current=-58
% HARDCODED[table: cf_table_baseline.tex, row: "Total Rewards", col: 1] current=-57

\begin{table}
    \caption{Counterfactual Effects on Prices and Shares}
    \label{tab:cf-effects}
    \begin{center}
        \small{
        \input{../output/tables/cf_table_baseline}
        }
    \end{center}
    \tablenote{Bootstrap standard errors are in parentheses.
    The "cap fees" scenario caps credit card merchant fees at 120 bp.
    The "uncap debit" scenario raises the cap on debit card merchant fees by $25$ bps.
    Monopoly refers to merging all three networks.
    Dual routing increases the multi-homing parameter by the equivalent of 20 bps in rewards for the median consumer.
    CBDC introduces a zero-fee public payment option.
    Its own market share is the residual of the changes in the shares of the other payment methods.
    Dollar values are computed by normalizing to \$10 trillion of total consumer purchases \parencite{Nilson2020a}.}
\end{table}

\subsubsection{Welfare}
\label{subsubsec:counterfactual-welfare}

\begin{table}[t]
    \caption{Welfare Effects of Counterfactual Policies}
    \label{tab:welfare-decomp}
    \begin{center}
        \small{
        \input{../output/tables/welfare_table_baseline}
        }
    \end{center}
    \tablenote{Bootstrap standard errors in parentheses.
    Cap Fees caps credit card merchant fees at 120 bp.
    Uncap Debit raises the debit fee cap by 25 bps.
    Monopoly merges all three networks.
    Dual Routing increases the multi-homing parameter by the equivalent of 20 bps in rewards for the median consumer.
    CBDC introduces a zero-fee public payment option.
    Low (high) income consumers have log income at $-2$ ($+2$) standard deviations relative to the median.
    Dollar values are computed by normalizing to \$10 trillion of total consumer purchases \parencite{Nilson2020a}.}
\end{table}

\begin{table}[t]
    \caption{Welfare Decomposition by Channel: 120 Bps Cap}
    \label{tab:cap-credit-channel-decomp}
    \begin{center}
        \small{
        \input{../output/tables/cap_credit_channel_decomp_baseline}
        }
    \end{center}
    \tablenote{Bootstrap standard errors in parentheses.
    Network Effect freezes merchant retail prices and acceptance at baseline while networks re-optimize fees and rewards.
    Price Passthrough then allows merchants to re-optimize retail prices, holding acceptance fixed.
    Merchant Adoption allows merchants to adjust acceptance decisions.
    All values in \$Bn, normalized to \$10 trillion of total consumer purchases \parencite{Nilson2020a}.}
\end{table}

The gains from the credit fee cap stem from correcting excessive adoption, not from reducing market power.
I measure consumer welfare by compensating variation:
\begin{equation}
    \text{CS}=\int\E\left[\max_{w}\log V^w_i\left(y\right)\right]\times y\d F\left( y\right)\label{eq:cons-surplus-integral}
\end{equation}
The inner expectation computes surplus for households at a given income as a share of baseline income. The outer integral weights by baseline income to yield aggregate welfare.
% HARDCODED[table: welfare_table_baseline.tex, row: "Total", col: 1] current=27
Most gains accrue to consumers (\$28 billion), not merchants (\$0.4 billion).
% HARDCODED[table: welfare_table_baseline.tex, row: "Consumers", col: 1] current=28
% HARDCODED[table: welfare_table_baseline.tex, row: "Merchants", col: 1] current=0.4
Networks lose only \$1.5 billion---around 6\% of baseline profits.
% HARDCODED[table: welfare_table_baseline.tex, row: "Networks", col: 1] current=-1.5
% HARDCODED[table: baseline_network_profits.tex, row: "Baseline", col: 1] current=26; 1.5/26 ≈ 0.058
The total gain is roughly double the \$12 billion consumer welfare gain from the CARD Act \parencite{Agarwal.etal2015}, highlighting the importance of regulating networks in addition to issuers.
% CONSTANT: from Agarwal et al. (2015), external paper result

Table~\ref{tab:cap-credit-channel-decomp} isolates the welfare sources by holding different margins fixed.
The first row holds merchant prices and acceptance at their baseline while networks re-optimize fees and rewards and consumers adopt new payment methods.
The second row additionally allows merchant retail prices to adjust, holding acceptance fixed.
The third row allows acceptance to adjust as well.

Even though the credit fee cap targets merchants, merchants benefit little in equilibrium.
The first row shows that merchants benefit from lower fees, but the second row shows that they dissipate these profits by lowering retail prices to compete for consumers.

The harm to networks is also modest.
The first row shows that lower rewards reduce profits as networks lose consumers.
The third row shows that networks recover a substantial share when merchants increase card acceptance in response to lower fees.
Profits decline by only \$1.5 billion, or around 6\% of baseline profits.
% HARDCODED[table: appendix_cap_welfare_table_baseline.tex, row: "Networks", col: 1] current=-1.5
% HARDCODED[the profit decline divided by table: baseline_network_profits.tex, row 1] current=-1.5 / 26

Consumers gain \$28 billion even though Table~\ref{tab:cf-effects} shows fees and rewards falling by roughly equal amounts.
The reason is that the credit fee cap moderates excessive adoption of credit cards.
By revealed preference, the marginal credit card user is indifferent between credit and debit or cash.
Rewards compensate for non-pecuniary costs of credit card use, including budget control, debt aversion, and account-management hassle.
I refer to these costs as ``credit aversion'' (Online Appendix~\ref{subsec:oa-survey-consumer-pref}).
Lower rewards cause this marginal consumer to switch away from credit.
But while lost rewards are merely transfers, the credit aversion that switchers avoid is a real social benefit.
The full total welfare gain of reduced credit aversion materializes only across the first two rows of Table~\ref{tab:cap-credit-channel-decomp}.\footnote{
    The first row of Table~\ref{tab:cap-credit-channel-decomp} conflates two offsetting forces: a large gain from reduced credit aversion and a loss from reduced consumption when rewards fall but retail prices remain fixed.
    Positive retail markups mean that reducing total consumption increases deadweight loss.
    However, this welfare channel is reversed in the second row as merchant price cuts then expand consumption.
}

Credit card adoption is socially excessive because of price coherence.
Because merchants do not surcharge, each new cardholder raises retail prices for everyone else.
Consumers collectively prefer lower retail prices and lower credit card use, but individually face incentives to defect to earn cross-subsidies.
The credit fee cap corrects this prisoner's dilemma by reducing the rewards that drive excessive adoption.

\subsubsection{Distributional Effects}
\label{subsubsec:counterfactual-distributional}

Fee and reward cuts redistribute consumption toward lower-income consumers.
Consumer surplus rises by 48 bp of consumption for low-income consumers (log income at $-2$ SD),
% HARDCODED[table: welfare_table_baseline.tex, row: "Low Income", col: 1] current=48
37 bp for median-income consumers,
% HARDCODED[table: welfare_table_baseline.tex, row: "Median Income", col: 1] current=37
and 15 bp for high-income consumers (log income at $+2$ SD).
% HARDCODED[table: welfare_table_baseline.tex, row: "High Income", col: 1] current=15
Lower retail prices benefit everyone, while reward reductions fall disproportionately on higher-income consumers who are more likely to use credit.
All groups gain on net because reduced credit aversion is spread across the income distribution.
To the extent that one believes revealed preference does not apply in payment markets, the average gain in consumer surplus can be subtracted from each group.
Such a correction would still leave gains for low- and median-income consumers, but losses for high-income consumers.

\subsection{Repealing the Durbin Amendment}
\label{subsec:repeal}

The credit fee cap raises welfare by correcting adoption distortions, not by reducing market power.
If they did, capping fees on any payment instrument should help.
The Durbin Amendment provides a direct test, since it caps debit but not credit card fees.
I study repealing the Durbin Amendment by raising the cap on debit card fees by $25$ bps, matching the share of the 75 bp interchange decline that flowed through to rewards \parencite{Hayashi2012}.\footnote{
    Durbin reduced debit interchange by roughly 75 bps, but not all of this flowed through to debit rewards.
    \label{foot:durbin-icg-vs-rewards}
    Banks also absorbed the loss through higher checking account fees \parencite{Kay.etal2018, Mukharlyamov.Sarin2025}.
    I ignore the increase in checking account fees since they do not influence the relative costs and benefits of different payment methods.
}
% CONSTANT: policy parameter, matches model scenario definition; empirical estimate from Mukharlyamov-Sarin (2025)

Repealing the Durbin Amendment moderates the rewards race.
If the debit fee cap were lifted by 25 bps, merchant fees would rise and networks would increase debit rewards.
Consumers would switch to debit, especially reward-sensitive ones.
As the marginal credit consumer becomes less reward-sensitive, networks pull back on rewards competition, triggering the see-saw principle and lowering credit merchant fees.
Credit rewards fall 6 bps and credit fees fall 3.3 bps as credit networks pull back.
% HARDCODED[table: cf_table_baseline.tex, row: "Credit Rewards", col: 3] current=-6
% HARDCODED[table: cf_table_baseline.tex, row: "Credit Fees", col: 3] current=-3.3
Net of the debit increase, total merchant fees rise by \$4.4 billion but rewards rise by only \$0.5 billion, reflecting substitution toward lower-reward debit.
% HARDCODED[table: cf_table_baseline.tex, row: "Total Fees", col: 3] current=4.4
% HARDCODED[table: cf_table_baseline.tex, row: "Total Rewards", col: 3] current=0.5

Repealing the Durbin Amendment generates large welfare gains.
Consumers gain \$4.3 billion as higher debit rewards draw consumers away from credit cards.
% HARDCODED[table: welfare_table_baseline.tex, row: "Consumers", col: 3] current=4.3
Total welfare rises by \$7 billion.
% HARDCODED[table: welfare_table_baseline.tex, row: "Total", col: 3] current=7
These gains are progressive: low-income consumers gain 7 bps of consumption, versus 2.4 bps for high-income consumers.
% HARDCODED[table: welfare_table_baseline.tex, row: "Low Income", col: 3] current=7
% HARDCODED[table: welfare_table_baseline.tex, row: "High Income", col: 3] current=2.4
The current regime, capping debit but not credit, is worse than both laissez-faire and capping both.
Uniform regulation benefits consumers \parencite{Rochet.Tirole2011}. Asymmetric regulation does not.

\subsection{Welfare Effects of Reducing Network Competition}
\label{subsec:competition}

One might expect that increasing competition among networks would substitute for fee regulation.
In payment markets, the opposite holds.
Competition fuels the rewards arms race that amplifies distortions from price coherence.
I test this by modeling a merger to monopoly, with fees and rewards set to maximize industry profits.

Without competitive pressure to fund rewards, a monopolist cuts credit card rewards by 90 bps.
% HARDCODED[table: cf_table_baseline.tex, row: "Credit Rewards", col: 6] current=-90
Because rewards are the primary reason consumers carry credit cards, spending on credit collapses. The credit share falls 23 pp.
% HARDCODED[table: cf_table_baseline.tex, row: "Credit", col: 6] current=-23
The merger reduces aggregate merchant fees by \$40 billion despite a 36 bps rise in per-transaction fees, because far fewer transactions clear on high-fee credit cards.
% HARDCODED[table: cf_table_baseline.tex, row: "Credit Fees", col: 6] current=36
% HARDCODED[table: cf_table_baseline.tex, row: "Total Fees", col: 6] current=-40
% HARDCODED[table: cf_table_baseline.tex, row: "Total Rewards", col: 6] current=-66
Unlike in \textcite{Armstrong2006}, where competition raises per-transaction merchant fees, here it raises the total fee burden by shifting payment composition toward high-fee credit cards.

The merger increases total welfare by \$15 billion because eliminating rewards competition more than offsets the consumer losses.
% HARDCODED[table: welfare_table_baseline.tex, row: "Total", col: 6] current=15
Consumers lose \$11 billion (though imprecisely estimated)
% HARDCODED[table: welfare_table_baseline.tex, row: "Consumers", col: 6] current=-11
while networks gain \$31 billion.
% HARDCODED[table: welfare_table_baseline.tex, row: "Networks", col: 6] current=31
Lower credit card use reduces retail prices for all consumers, but lower rewards mainly harm high-income cardholders.
Low-income consumers gain 13 bps in consumption
% HARDCODED[table: welfare_table_baseline.tex, row: "Low Income", col: 6] current=13
while high-income consumers lose 29 bps.
% HARDCODED[table: welfare_table_baseline.tex, row: "High Income", col: 6] current=-29
These results do not support mergers, since the consumer losses are large and imprecisely estimated.
They do show that rewards competition is so harmful that increases in market power can raise total welfare.

\subsection{Consumer Multi-Homing and Dual Routing}
\label{subsec:dual-routing}

Rather than regulating fees directly, policy can redirect the locus of network competition.
Routing requirements, like those in the Credit Card Competition Act, would force banks to issue cards that are compatible on multiple networks \parencite{Andriotis2022}.
This increases effective multi-homing rates and allows merchants to route payments to cheaper networks, shifting the locus of competition from the consumer side toward merchant fees \parencite{Teh.etal2022}.
I model this by raising the $\chi^{22}$ parameter, which captures the complementarity of holding two credit cards, thereby increasing consumer multi-homing.
For the median consumer, this is equivalent to a 20 bps reward increase for those bundles.
Credit card multihoming then rises by \scalarinput{dual_routing_cc_multihome_change_baseline} pp. from a baseline level of around \scalarinput{kilts_top_two_many_cc}\%.
% CONSTANT: 20 bps reward increase is policy parameter for this counterfactual scenario
Credit card fees fall 13 bps
% HARDCODED[table: cf_table_baseline.tex, row: "Credit Fees", col: 8] current=-13
and rewards fall 15 bps.
% HARDCODED[table: cf_table_baseline.tex, row: "Credit Rewards", col: 8] current=-15
Consumer welfare, excluding the mechanical impact of increasing $\chi^{22}$, still rises \$3.3 billion, and total welfare rises \$3.8 billion.\footnote{
% HARDCODED[table: welfare_table_baseline.tex, row: "Consumers", col: Dual Routing] current=3.3
% HARDCODED[table: welfare_table_baseline.tex, row: "Total", col: Dual Routing] current=3.8
Because dual routing directly modifies $\chi^{22}$, it mechanically raises consumer surplus even at unchanged prices and portfolios. I net this mechanical component---around \$4.5 billion---out of the reported consumer surplus so that the \$3.3 billion figure reflects only the equilibrium response. Online Appendix~\ref{subsec:oa-mechanical-dr} derives the adjustment.
To the extent that changing preferences complicates the interpretation of welfare, the effects on fees and rewards are more robust as they are a consequence of how merchant acceptance depends on consumer multi-homing across networks.}
% HARDCODED[table: mechanical_dual_routing_baseline.tex, row: Aggregate] current=4.5

The intuition for these effects follows from the merchant acceptance condition in Section \ref{subsubsec:model-merch-accept}.
A merchant's fee cost from accepting card $j$ depends on all consumers who use it, but the incremental sales gain comes only from single-homers who would otherwise pay cash.
When consumers single-home on one credit network, merchants must accept that network to capture those sales---what trial testimony in Ohio v.\ American Express called ``cardholder insistence'' \parencite{Conrath2014}.
When a consumer carries credit cards from multiple networks, the merchant can decline the expensive network's card without losing the sale---the consumer simply pays with the cheaper alternative.
The incremental sales gain from accepting the expensive network shrinks.
Networks must then compete on merchant fees rather than relying on consumer lock-in to extract rents from the merchant side.

Current proposed dual-routing legislation often requires that the secondary network not be a large established network.
This exclusion is unnecessary. 
The power of the dual-routing mechanism stems from increasing multi-homing, not from reducing the size of the largest networks.
The exclusion may even be counterproductive if merchants are unwilling to invest in accepting smaller networks.

\subsection{Central Bank Digital Currencies and Public Entry}
\label{subsec:cf-cbdc}
Proposals for central bank digital currencies (CBDC) and faster payments \parencite{Shin2021,FederalReserve2022,Berg.etal2024a} motivate the possibility that government entry could discipline network pricing.
I simulate a public debit network with MC debit's demand and cost characteristics that sets merchant fees at cost and offers zero consumer rewards, earning zero profit.
Because the entrant offers no consumer-side subsidies, its resulting market share is small.
Consumer welfare gains are roughly \$3.7 billion and total welfare gains are roughly \$1.6 billion, both smaller than repealing Durbin caps and smaller than the \$3.8 billion total welfare gain under dual routing.
% HARDCODED[table: welfare_table_baseline.tex, row: "Consumers", col: CBDC] current=3.7
% HARDCODED[table: welfare_table_baseline.tex, row: "Total", col: CBDC] current=1.6
% HARDCODED[table: welfare_table_baseline.tex, row: "Total", col: Dual Routing] current=3.8
These results depend on modeling the entrant as a debit product that does not substitute for credit at the point of sale.
The segmentation assumption (Section~\ref{subsec:model-assumptions}) limits the competitive pressure a new payment method can exert on incumbent merchant fees.
A public product that also displaced credit could generate larger inroads, but is beyond the scope of most of the current proposals for public options.
The estimated welfare gains do not include any fixed costs of setting up the public network.
Directly regulating fees or mandating multi-homing is thus likely more effective than introducing a public option.

\subsection{Discussion}
\label{subsec:cf-discussion}

Online Appendix \ref{sec:oa-alternative-specifications} re-estimates the model under alternative assumptions about consumer constraints, pass-through, and reward sensitivity.

\paragraph{Constraints vs Preferences}
\label{subsec:revealed-pref-approx}
Consumer payment choices in the model reflect preferences, not constraints.
Some consumers' choice to primarily use debit over credit is evidence of ``credit aversion.''
Survey evidence documents concrete motives, including budget control, debt aversion, and account-management hassle (Online Appendix \ref{subsec:oa-survey-consumer-pref}).

An alternative interpretation is that some consumers cannot obtain credit cards.
The presence of exogenous constraints does not affect the estimated welfare effects.
Online Appendix~\ref{subsec:oa-credit-constrained} uses credit score data from the DCPC to estimate a model in which some consumers are credit-constrained, with the extent of the constraint varying with income.
The alternative specification predicts smaller credit aversion for the median consumer (lower $\Xi$ for credit and a weaker income gradient $\beta_y^L$).
Because constrained consumers do not respond on the credit margin, the unconstrained subset must be more reward-sensitive to rationalize the same Durbin moment.
The estimated consumer and total welfare results are similar to the baseline model.


\paragraph{Pass-through}
\label{subsubsec:passthrough-theory}
The CES functional form implies full pass-through from merchant fees to retail prices.
One reason full pass-through is a reasonable baseline for interchange is because interchange is an aggregate shock to all merchants \parencite{Amiti.etal2019}.
I also lack the required merchant-level interchange data matched to retail prices to test pass-through.
Online Appendix \ref{subsec:oa-no-passthrough} compares the full pass-through baseline against zero pass-through, in which retail prices are unresponsive to merchant fees in either direction.
The alternative shifts the consumer-merchant welfare split but yields similar equilibrium fees and rewards.
% HARDCODED[table: welfare_table_no_passthrough.tex vs welfare_table_baseline.tex, row: Total, col: Cap Fees, (27-17)/27 ≈ 1/3] current=one-third

\paragraph{Reward Sensitivity}
\label{subsubsec:reward-sensitivity}
A more conservative specification targets only half the observed Durbin decline, attributing the remainder to non-price steering.
Online Appendix~\ref{subsec:oa-debit-robustness} re-estimates the model against this lower-bound target.
The halved target produces negative marginal costs for Visa, which conflict with accounting data on issuer, acquirer, and network costs.
The credit fee cap remains welfare-improving in this specification.
The monopoly counterfactual is the most sensitive to this calibration: its welfare gain reverses sign under the halved target.
If non-rewards steering imposes real costs without direct consumer benefits, these estimates understate the true losses from high merchant fees.

\paragraph{Choice of Credit Fee Cap}
The 120 bp cap keeps the counterfactual close to observed fee levels, avoiding the large out-of-sample extrapolation required by more aggressive caps.
Online Appendix~\ref{subsec:oa-cost-of-cash-cap} compares this cap to a 30 bps cap modeled after the tourist test of \textcite{Rochet.Tirole2011} and to the Ramsey planner's solution.
The 120 bps cap achieves roughly 93\% and 77\% of the respective welfare gains with less out-of-sample extrapolation.
% HARDCODED[table: appendix_cap_welfare_table_baseline.tex, row: "Total", cols: Cap Credit=27, Tourist Test=29, Ramsey Planner=35] current=93% (27/29), 77% (27/35)

\end{document}
